% GENERATED FILE. Edit canonical lessons, then run npm run generate:books.
% canonical-source-hash: fnv1a64:991eb6d7c8df8acd
% canonical-lessons: ES-C756-delantero, ES-C756-marcar, ES-C756-marcador, ES-C756-arbitro, ES-C756-empatar

\chapter{El partido}
\label{ch:es-el-partido}
\hlchaptermodality{756}

\begin{chapteropening}
\textbf{By the end of this chapter:} \emph{I can follow a football match: the striker, scoring, the score, the referee and a draw.}

Complete the last lesson of chapter 756: I can follow a football match: the striker, scoring, the score, the referee and a draw.
\end{chapteropening}

\section[el delantero]{el delantero — the forward, the striker}

\label{lesson:ES-C756-delantero}



\begin{quote}
Before the new one: say the Spanish for to get injured, then the Spanish for the muscle.
\end{quote}

\subsection*{You'll want to know: el delantero}
\textbf{el delantero} — \textquotedblleft{}the forward, the striker\textquotedblright{}.

\begin{quote}
Nuestro delantero es muy rápido.
\end{quote}

\begin{quote}
\emph{Our striker is very fast.}
\end{quote}

As an adjective: \emph{la parte delantera} — the front.

\begin{cousinweb}
\emph{delantero} is built on \emph{delante} \textquotedblleft{}in front\textquotedblright{}, which comes from Old Spanish \emph{denante} (EtymDB), from Latin \emph{dē}, \emph{in} and \emph{ante} \textquotedblleft{}before\textquotedblright{}.
\end{cousinweb}

\begin{grammarlens}[title={What you've built}]
One more word for telling, explaining and describing at B1.
\end{grammarlens}

\subsection*{Your turn}
\begin{itemize}
\raggedright
  \item \emph{Say it:} \emph{el delantero}
  \item \emph{Say it:} \emph{el delantero}, once more
  \item \emph{Say it:} a sentence of your own with \emph{el delantero}
  \item \emph{Recall:} say the Spanish for to get injured, then the Spanish for the muscle, then say \emph{el delantero} again
\end{itemize}

\begin{tcolorbox}[breakable,colback=teal!4,colframe=teal!35!black,title={Before you move on}]
What does el delantero mean? (\textquotedblleft{}The forward, the striker\textquotedblright{}.) Say it once more.
\end{tcolorbox}
\section[marcar]{marcar — to score; to mark; to dial}

\label{lesson:ES-C756-marcar}



\begin{quote}
Before the new one: say the Spanish for the muscle, then the Spanish for the forward.
\end{quote}

\subsection*{You'll want to know: marcar}
\textbf{marcar} — \textquotedblleft{}to score; to mark; to dial\textquotedblright{}.

\begin{quote}
Marcó dos goles en la final.
\end{quote}

\begin{quote}
\emph{He scored two goals in the final.}
\end{quote}

Also to dial a number: \emph{marcar el número}. Spelling: \emph{marqué}.

\begin{cousinweb}
\emph{marcar} comes from Italian \emph{marcare} \textquotedblleft{}to mark\textquotedblright{} (EtymDB), which is built on Germanic \emph{marka} \textquotedblleft{}boundary, sign\textquotedblright{}.
\end{cousinweb}

\begin{grammarlens}[title={What you've built}]
One more word for telling, explaining and describing at B1.
\end{grammarlens}

\subsection*{Your turn}
\begin{itemize}
\raggedright
  \item \emph{Say it:} \emph{marcar}
  \item \emph{Say it:} \emph{marcar}, once more
  \item \emph{Say it:} a sentence of your own with \emph{marcar}
  \item \emph{Recall:} say the Spanish for the muscle, then the Spanish for the forward, then say \emph{marcar} again
\end{itemize}

\begin{tcolorbox}[breakable,colback=teal!4,colframe=teal!35!black,title={Before you move on}]
What does marcar mean? (\textquotedblleft{}To score; to mark; to dial\textquotedblright{}.) Say it once more.
\end{tcolorbox}
\section[el marcador]{el marcador — the scoreboard, the score}

\label{lesson:ES-C756-marcador}



\begin{quote}
Before the new one: say the Spanish for the forward, then the Spanish for to score.
\end{quote}

\subsection*{You'll want to know: el marcador}
\textbf{el marcador} — \textquotedblleft{}the scoreboard, the score\textquotedblright{}.

\begin{quote}
El marcador final fue dos a uno.
\end{quote}

\begin{quote}
\emph{The final score was two-one.}
\end{quote}

In Latin America \emph{un marcador} is also a marker pen.

\begin{cousinweb}
\emph{marcador} is built on \emph{marcar} with the ending \emph{-dor}; \emph{marcar} comes from Italian \emph{marcare}, built on Germanic \emph{marka} \textquotedblleft{}boundary, sign\textquotedblright{}.
\end{cousinweb}

\begin{grammarlens}[title={What you've built}]
One more word for telling, explaining and describing at B1.
\end{grammarlens}

\subsection*{Your turn}
\begin{itemize}
\raggedright
  \item \emph{Say it:} \emph{el marcador}
  \item \emph{Say it:} \emph{el marcador}, once more
  \item \emph{Say it:} a sentence of your own with \emph{el marcador}
  \item \emph{Recall:} say the Spanish for the forward, then the Spanish for to score, then say \emph{el marcador} again
\end{itemize}

\begin{tcolorbox}[breakable,colback=teal!4,colframe=teal!35!black,title={Before you move on}]
What does el marcador mean? (\textquotedblleft{}The scoreboard, the score\textquotedblright{}.) Say it once more.
\end{tcolorbox}
\section[el árbitro]{el árbitro — the referee}

\label{lesson:ES-C756-arbitro}



\begin{quote}
Before the new one: say the Spanish for to score, then the Spanish for the scoreboard.
\end{quote}

\subsection*{You'll want to know: el árbitro}
\textbf{el árbitro} — \textquotedblleft{}the referee\textquotedblright{}.

\begin{quote}
El árbitro no vio la falta.
\end{quote}

\begin{quote}
\emph{The referee did not spot the foul.}
\end{quote}

\emph{La árbitra} for a woman.

\begin{cousinweb}
\emph{árbitro} comes from Latin \emph{arbiter} \textquotedblleft{}judge, umpire\textquotedblright{} (EtymDB).
\end{cousinweb}

\begin{grammarlens}[title={What you've built}]
One more word for telling, explaining and describing at B1.
\end{grammarlens}

\subsection*{Your turn}
\begin{itemize}
\raggedright
  \item \emph{Say it:} \emph{el árbitro}
  \item \emph{Say it:} \emph{el árbitro}, once more
  \item \emph{Say it:} a sentence of your own with \emph{el árbitro}
  \item \emph{Recall:} say the Spanish for to score, then the Spanish for the scoreboard, then say \emph{el árbitro} again
\end{itemize}

\begin{tcolorbox}[breakable,colback=teal!4,colframe=teal!35!black,title={Before you move on}]
What does el árbitro mean? (\textquotedblleft{}The referee\textquotedblright{}.) Say it once more.
\end{tcolorbox}
\section[empatar]{empatar — to draw, to tie}

\label{lesson:ES-C756-empatar}



\begin{quote}
Before the new one: say the Spanish for the scoreboard, then the Spanish for the referee.
\end{quote}

\subsection*{You'll want to know: empatar}
\textbf{empatar} — \textquotedblleft{}to draw, to tie\textquotedblright{}.

\begin{quote}
Empatamos uno a uno.
\end{quote}

\begin{quote}
\emph{We drew one all.}
\end{quote}

\emph{El empate} — the draw.

\begin{cousinweb}
\emph{empatar} was borrowed from Italian \emph{impattare} \textquotedblleft{}to draw\textquotedblright{}, built on \emph{patta} \textquotedblleft{}a draw, a tie\textquotedblright{} (EtymDB). Where \emph{patta} comes from is not settled: EtymDB lists both a Lombardic word and Latin \emph{pactum} \textquotedblleft{}agreement\textquotedblright{}.
\end{cousinweb}

\begin{grammarlens}[title={What you've built}]
One more word for telling, explaining and describing at B1.
\end{grammarlens}

\subsection*{Your turn}
\begin{itemize}
\raggedright
  \item \emph{Say it:} \emph{empatar}
  \item \emph{Say it:} \emph{empatar}, once more
  \item \emph{Say it:} a sentence of your own with \emph{empatar}
  \item \emph{Recall:} say the Spanish for the scoreboard, then the Spanish for the referee, then say \emph{empatar} again
\end{itemize}

\begin{tcolorbox}[breakable,colback=teal!4,colframe=teal!35!black,title={Before you move on}]
What does empatar mean? (\textquotedblleft{}To draw, to tie\textquotedblright{}.) Say it once more.
\end{tcolorbox}
